% =====================================================================
\section{Introduction}\label{sec:introduction}
% =====================================================================
In an automated pick-and-pass warehouse a conveyor carries each order box past fixed picking stations, and the box diverts to every station that stores one of its products. The submitted companion paper \citep{belul_companion} poses correlated storage assignment for such systems with total order-station visits as the objective, under hard per-station limits on slots and workload, and applies it to an operating warehouse, Company~A. The workload limit exists because a layout that concentrates co-ordered products to save visits can route traffic to a few stations: a pick-and-pass line behaves as a network of queues \citep{yu2008pickandpass}, and a station whose arrival rate approaches its service rate blocks the main line \citep{kim2020}.

A layout is fitted to orders that have already been observed and is then operated on orders that come later. What the workload limit has to protect is therefore the distribution of work across stations on those later orders. The companion states the limit as a budget in workload units: each station's budget is 110\% of the complete legacy load it carried \citep{belul_companion}. Such a budget is tied to the volume of the horizon on which it was calibrated, whereas the volume of a future horizon is not known when the layout is chosen. This article studies a rule stated in shares instead: after re-slotting, each station should keep its share of future workload within a fixed number of percentage points of the share it carried historically under the layout the site operates. % AUTHOR-CONFIRM: whether Company A states its operational workload policy in these terms (AUTHOR_QUESTIONS.md, item 1).
A share rule needs no absolute workload level, since multiplying every product's workload by a common positive factor leaves every share unchanged. It does not make the horizon irrelevant, because a different horizon brings a different product mix.

The companion already evaluates its layouts out of sample: on a dated extract, a layout fitted to the first weeks and scored on the following weeks removes 6.6\% and 7.4\% of station visits, against 13.7\% in sample \citep{belul_companion}. Scoring layouts on later orders is thus already part of the companion's design. What the companion does not address is a share rule on later orders, and the choice between the ways of protecting such a rule that a planner can build into the optimisation. Three features of the problem shape that choice. First, the visit objective gives an optimiser no reason to leave part of a workload allowance unused, so a layout that satisfies a share rule on history may sit at the edge of the band at some station, where a small change in the product mix carries it outside. Second, the rule is two-sided: shares sum to one, so caps alone would let one station lose a large part of its work while the others absorb it, and a floor has to be protected as explicitly as a cap. Third, no static layout can keep a nontrivial band for every conceivable future, since a future whose work concentrates on the products of one station drives that station's share towards one; any protection therefore rests on an assumption about which futures matter.

Two standard responses are available. One reserves part of the band during optimisation and scores the layout at the full band, a device related to the protection levels of budgeted robust optimisation \citep{bertsimas2004price}. The other requires the band to hold for every product mix in a set built from historical variation, the set-inclusive form of robust optimisation \citep{soyster1973convex}. Here that set is the convex hull of the product mixes of historical order blocks, extended by activation protection, an allowance of workload mass on products that had no historical workload. The closest storage-assignment work states station balance on average demand and finds it exceeded under simulated demand variation \citep{winkelmann2025integrated}; Section~\ref{sec:related} positions the article against that work and the wider literature. Whether a reserved margin, historical scenarios, or both keep a visit-minimising layout inside a share band on later orders is an empirical question for a given site, and it is the question of this article.

We extend the companion's model to a share rule and compare the two responses on Company~A's order stream. The visit objective and the assignment structure are retained, and the per-station budget rows are replaced by rows that keep each station's share inside a two-sided absolute band around its historical target (Section~\ref{sec:model}). Workload counts product-order pairs rather than the companion's load normalised by processing rate, so the band is stated in shares of order lines, not in station utilisation. The rows are stated either over pooled history alone or over historical block scenarios with activation protection, and they are imposed either at the full band or at a band tightened by a reserved margin. An information contract fixes what is known when a layout is chosen: conditional on a snapshot of the catalogue and the incumbent layout, targets, training workloads and scenarios come from the orders before a deployment origin; the layout is frozen and then scored on the next $n$ complete orders (Section~\ref{sec:protocol}). Exploratory campaigns at an earlier origin shaped the rule and its parameters. A later origin then carried a predeclared two-by-two comparison of four model configurations, called arms, that crosses the reserved margin with scenario-and-activation protection; it was solved with three optimiser seeds and scored on orders that no earlier solve, score, or drift or dispersion analysis of the study had used.

% Contribution statement, revised in review round 1: held-out comparison first, formulation second, diagnostics third.
The article is a bounded comparative study, and its contribution has three parts. The first is held-out evidence on the two responses. Both arms that reserved a margin satisfied the band on all three optimiser seeds and neither arm without one did on any seed; in additional descriptive analyses, the arms with historical scenarios and activation kept their largest station deviations from target lower than the matching arms without them, at higher mean visits per order. The second is the share-policy formulation and information contract that make this comparison well posed: frozen targets, a closed catalogue in which frozen products keep their stations while their demand still counts, and a finite robust counterpart whose feasibility is conditional on the future product mix lying in the declared set. Its ingredients are standard (Section~\ref{subsec:rw_position}). The third is a set of diagnostics that bound the reading: the later product mix lay outside the set of mixes each returned layout was protected against, a sufficient condition on product-mix distance for the margin did not hold on any analysed window, and no validated rule for choosing the margin was obtained.

Section~\ref{sec:related} positions the study against the companion and the closest literature. Section~\ref{sec:model} states the information contract and the formulation. Section~\ref{sec:protocol} describes the data, the protocol, the staged design, the metrics and the scope of the evidence. Section~\ref{sec:results} reports the held-out comparison, then the exploratory context and additional descriptive analyses. Section~\ref{sec:discussion} discusses the findings and their limitations, and Section~\ref{sec:conclusion} concludes.
